\documentclass{report}

% ======================================
% PREÁMBULO: PAQUETES (Válido para todo el documento)
% ======================================
\usepackage[utf8]{inputenc}
\usepackage{amsmath}
\usepackage{amsfonts}
\usepackage{amssymb}
\usepackage{graphicx}
\usepackage{xcolor} % Para usar colores
\usepackage[margin=1in]{geometry} % Para controlar los márgenes
\usepackage{fancyhdr} % Opcional: para el logo en el encabezado

\pagestyle{fancy} 
\fancyhf{} % Limpia todos los encabezados y pies existentes

% 2. Configurar el contenido del encabezado
% \leftmark muestra el nombre del Capítulo
% \rightmark muestra el nombre de la Sección
\fancyhead[L]{\nouppercase{\leftmark}} % Nombre del Capítulo (Izquierda)
%\fancyhead[R]{\nouppercase{\rightmark}} % Nombre de la Sección (Derecha)

% 3. Configurar el pie de página
\fancyfoot[C]{\thepage} % Número de página centrado

% 4. Regla de separación
\renewcommand{\headrulewidth}{0.4pt} 

% 5. CORRECCIÓN CRÍTICA: Redefinir el estilo 'plain'
% El estilo 'plain' se usa por defecto en la primera página de cada \chapter y anula a 'fancy'.
% Debemos redefinirlo para que use nuestra configuración o simplemente limpie el encabezado.
\fancypagestyle{plain}{
	\fancyhead{} % Deja el encabezado de la primera página del capítulo vacío
	\fancyfoot[C]{\thepage} % Mantiene el número de página
	\renewcommand{\headrulewidth}{0pt} % Sin línea de separación en esa página
}



\usepackage{hyperref} 
\usepackage{listings} 
\usepackage[spanish]{babel} 
\usepackage{csquotes}
\usepackage[T1]{fontenc} 
\usepackage{placeins}
\definecolor{codegray}{rgb}{0.5,0.5,0.5}
\lstset{
	language=C, % Especifica el lenguaje
	basicstyle=\footnotesize\ttfamily, % Fuente y tamaño
	numbers=left, % Muestra números de línea a la izquierda
	numberstyle=\tiny\color{codegray}, % Estilo de los números de línea
	stepnumber=1, % Incremento de los números de línea
	showspaces=false, % No muestra espacios
	showtabs=false, % No muestra tabulaciones
	frame=single, % Dibuja un marco alrededor del código
	tabsize=2, % Tamaño de la tabulación
	captionpos=b, % Posición del título (bottom)
	breaklines=true, % Permite que las líneas largas se rompan
	breakatwhitespace=false, % Solo rompe en espacios
	keywordstyle=\color{blue}\bfseries, % Estilo para palabras clave
	identifierstyle=\color{black}, % Estilo para identificadores
	commentstyle=\color{teal}, % Estilo para comentarios
	stringstyle=\color{orange}, % Estilo para cadenas
	rulecolor=\color{black}, % Color del marco
	backgroundcolor=\color{gray!10} % Color de fondo (opcional, requiere \usepackage{xcolor})
}
\usepackage{multirow}
\graphicspath{{images/}}
% Estilos y configuraciones de listings, hyperref, etc.
\lstdefinestyle{myCode}{
	basicstyle=\ttfamily\footnotesize,
	keywordstyle=\color{blue},
	commentstyle=\color{green!50!black},
	stringstyle=\color{red},
	breaklines=true,
	showspaces=false,
	showtabs=false,
	tabsize=4
}
\renewcommand{\lstlistingname}{Código}
\lstset{style=myCode}

\usepackage{url}

\hypersetup{
	colorlinks=true,
	citecolor=black,
	filecolor=black,
	linkcolor=black,
	urlcolor=blue
}


% Puedes usar \include en lugar de \input para que cada capítulo empiece en una nueva página.
% \include es más eficiente si tu documento tiene muchos capítulos y quieres que cada uno fuerce una nueva página.
% \input simplemente inserta el contenido en el punto donde se llama.
\newcommand{\incluirCapitulo}[1]{\include{#1}} 
% Si prefieres \input (sin forzar nueva página, solo inserta el texto):
% \newcommand{\incluirCapitulo}[1]{\input{#1}}

% Definición del color para la portada
\definecolor{franjagris}{rgb}{0.85, 0.85, 0.85} 

\begin{document}
	
	% ======================================
	% PORTADA
	% ======================================
	\thispagestyle{empty} % Asegura que esta página no tenga números ni encabezados/pies de página
	\begin{titlepage}
		
		% 1. Logo/Incio de imagen (alineado a la derecha en la parte superior)
	\begin{flushright}
		\vspace*{0.5cm} % Espacio desde el borde superior
		
		% Coloca \includegraphics directamente aquí, sin el entorno figure
		\includegraphics[width=0.2\linewidth]{images/Logo}
		
		\vspace{0.2cm} % Espacio entre la imagen y el texto
	\end{flushright}
		
		\vfill % Empuja el contenido restante hacia abajo
		
		% 2. Contenido principal (alineado a la izquierda, cerca de la parte inferior)
		% Franja gris con el título
\colorbox{franjagris}{
	\begin{minipage}{\textwidth}
		\raggedleft % CAMBIO AQUÍ: Alinea el texto a la derecha
		\vspace{0.8cm} % Espacio superior dentro de la caja
		\Large \textbf{WaveLink32:} \\
		\vspace{0.2cm}
		\Huge \textbf{Generación de código para la STM32 usando Embedded Coder} \\
		\vspace{0.8cm} % Espacio inferior dentro de la caja
	\end{minipage}
}
		
		\vspace{0.5cm} % Espacio entre la franja y el texto del pie
		
	\raggedleft 
	Manual que describe como generar código C para los periféricos que maneja la biblioteca del \\ STM32 Embedded Processor de Simulink.
	\end{titlepage}
	
	
	% ======================================
	% INICIO DEL CONTENIDO / CAPÍTULOS
	% ======================================
	
	% *IMPORTANTE*: 
	% A partir de aquí, define el estilo de página para el resto del documento.
	% Por ejemplo, si quieres el número de página simple:
	%\pagestyle{plain} 
	
	% Tabla de Contenidos (opcional)
	\tableofcontents 
	\newpage
	
	% LLAMADAS A LOS CAPÍTULOS
	% Simplemente comenta o descomenta la línea de \incluirCapitulo para incluir o no el capítulo.
	\incluirCapitulo{Introduccion}
	\incluirCapitulo{IntrofamiliasStm32}
	\incluirCapitulo{Instalacion}
	\incluirCapitulo{EmbeddedExplanation}
	\incluirCapitulo{Conceptos}
	\incluirCapitulo{Macros}
	\incluirCapitulo{IntroplantillaCubeMX}
	\incluirCapitulo{ModeloPrueba}
	%\incluirCapitulo{TLK1}
	\incluirCapitulo{escrituraDigital}
	\incluirCapitulo{lecturaDigital}
	%\incluirCapitulo{TLK123}
	\incluirCapitulo{TLK124}
	\incluirCapitulo{adcManyPins}
	%\incluirCapitulo{TLK2}
	
	
	%\incluirCapitulo{portadavalid}
%	\incluirCapitulo{portadademo}
%	\incluirCapitulo{demo}
	% \incluirCapitulo{capitulo3} % Este capítulo está comentado y NO se incluirá en la compilación.
	
	
\end{document}